% ===== EXHIBIT: the governed dial (Appendix E per owner decision O-J3) =====
% Source: replication/run_admission_exhibits.py on cma_data snapshot 2026q2_custom.

\begin{figure}[H]
	\captionof{figure}{The Governed Dial: Cap 3 Sets the Admissible Admission Range}\label{fig:governed_dial}\vspace*{-0.5\baselineskip}
	\begin{center}
		\includegraphics[width=0.90\linewidth]{figures/governed_dial.PNG}
	\end{center}
	\vspace*{-1.0\baselineskip}
	{\footnotesize Notes: 2026-Q2 production cut, 18 assets, USD, Balanced with Alternatives mandate. For each admitted sleeve we sweep its admission weight $w_i$ over $[0,1]$ with the other sleeves held at the production policy, re-solve the mandate, and record the largest $w_i$ that satisfies each governance constraint. Nested bars are the Cap~3 portfolio budget $\boldsymbol\alpha_{adm}^{\intercal}\boldsymbol D^{-1}\boldsymbol\alpha_{adm} \le \kappa \cdot SR^2_{MATFCMA}$ at the three budgets, with the OTHER sleeves held at their Cap~3-projected weights $\theta(\kappa) w_j$, so a shorter bar is a tighter budget. Held instead at the production policy, the other sleeves alone exhaust the budget at every $\kappa$ shown and no admission is admissible: the sleeve-level restatement of the Cap~3 audit result. Orange diamonds mark the production admission. The benchmark stress floor of Section~5.3, $\text{stress total} \ge \mu_{bm} - 1.0\sigma_{bm}$, evaluates to -1.62\% of expected total return. the one-sigma benchmark stress floor never binds across the sweep, so Cap~3 is the operative constraint on every sleeve.}
\end{figure}
